\chapter{Integration and Outlook}
\label{ch:integration}

\section{Complementarity of the Two Frameworks}
\label{sec:integration_compare}

Chapters~\ref{ch:heine} and~\ref{ch:dieckow} attack the same problem --- inferring
species interactions from compositional time series under the Extended Hamilton
Principle --- from opposite ends of the controlled/clinical spectrum. They are
deliberately complementary rather than redundant (Table~\ref{tab:complementarity}).

\begin{table}[htbp]
  \centering\small
  \begin{tabular}{@{}lll@{}}
    \toprule
    \textbf{Aspect} & \textbf{Ch.~\ref{ch:heine} (Heine, in-vitro)} & \textbf{Ch.~\ref{ch:dieckow} (Dieckow, in-vivo)}\\
    \midrule
    Data           & in-vitro, controlled & in-vivo, clinical\\
    Dimension      & $5$ species          & $10$ class-level guilds\\
    Model          & Hamilton $\bphi$--$\bpsi$ (viability) & Hamilton replicator ($\bphi$ only)\\
    Inference      & Bayesian TMCMC (full posterior) & L-BFGS-B point estimate $+$ LOO\\
    Regulariser    & iteratively narrowed priors & metabolic FBA \emph{sign} prior\\
    Acceleration   & GPU \texttt{vmap} ($\sim$$200\times$) & JAX-JIT autodiff\\
    Validation     & $4$ conditions $+$ held-out pH & $10$-fold LOO $+$ Joshi external\\
    \bottomrule
  \end{tabular}
  \caption[Complementarity of the in-vitro and in-vivo frameworks]{The two
  inference frameworks are complementary: controlled identifiability and full
  posteriors in vitro; clinical relevance and metabolic-prior identifiability in
  vivo.}
  \label{tab:complementarity}
\end{table}

The in-vitro setting (Chapter~\ref{ch:heine}) trades realism for control. With
defined strains, replicated conditions, and measured viability, it can afford a
\emph{full Bayesian posterior} over a $15$-parameter symmetric matrix and a
$\bphi$--$\bpsi$ model that resolves living from total biomass --- at the cost of
applying only to a five-species laboratory consortium. The data are rich enough
that the regulariser need only be an iteratively narrowed prior; the posterior
itself recovers biological sparsity without metabolic input.

The in-vivo setting (Chapter~\ref{ch:dieckow}) inverts every term of that trade.
Ten guilds, three time points, and ten patients leave the interaction matrix
underdetermined, so a point estimate replaces the posterior and an external
\emph{metabolic sign prior} supplies the identifiability the data cannot. What the
clinical setting loses in statistical luxury it gains in reach: out-of-sample LOO
generalisation and transfer to an independent peri-implant cohort.

Crucially, the two share a backbone. Both are replicator dynamics on the simplex
with a symmetric interaction matrix derived from the same variational principle
(Chapter~\ref{ch:theory}), and both find that interaction \emph{sign} is the
robust, transferable quantity --- recovered emergently from rich in-vitro data and
imposed as a prior on sparse in-vivo data. Read together they form a single
pipeline: controlled in-vitro estimation calibrates the model and validates the
sign structure that the in-vivo prior then exploits where data are scarce, and the
spatial extension of this chapter carries the same interactions into depth-resolved
prediction.
A structural echo of this complementarity is visible in the COMSTAT biofilm metrics:
the roughness coefficient $R_a$ of \textit{P.~gingivalis} is significantly higher
in the dysbiotic-HOBIC condition than in the commensal one (Mann--Whitney
$p < 0.001$), which mirrors the Hamilton posterior showing a stronger
Fn$\to$Pg interaction in DH ($\mathrm{P}(a_{Fn\to Pg}>0)=1.00$) versus CH
($\mathrm{P}=0.82$) --- a suggestive co-variation between spatial surface roughness
and inferred network interaction strength, consistent with the picture that
stronger Fn facilitation allows Pg to form more structurally complex colonies.

A deeper complementarity concerns what each framework \emph{measures}.
The Hamilton ODE equilibrium (300 posterior samples, uniform initialisation)
converges to Pg~$+$~Fn dominance
($\phi_\mathrm{Pg}^\mathrm{CH}\!\approx 0.49$; $\phi_\mathrm{Fn+Pg}^\mathrm{DH}\!\approx 1.0$),
yet COMSTAT shows near-uniform substrate coverage across all five species
($0.19$--$0.26$; $r = 0.07$ CH, $r = 0.52$ DH with the ODE).
The discrepancy is resolved by vertical stratification: competitive exclusion in the
mean-field model corresponds to depth-niche partitioning in the real biofilm,
which COMSTAT's 2-D projection cannot detect but FISH can (Section~\ref{sec:pde_fish}).
Figure~\ref{fig:hamilton_vs_comstat} shows the ODE--COMSTAT correlation explicitly.

\begin{figure}[htbp]
  \centering
  \includegraphics[width=0.74\textwidth]{fig_hamilton_vs_comstat.pdf}
  \caption[Hamilton ODE equilibrium vs COMSTAT substrate coverage]{Hamilton ODE
  equilibrium composition (posterior median, $n=300$ samples) versus COMSTAT
  substrate coverage fraction for all five species in CH and DH.
  The low correlation ($r=0.07$ CH, $r=0.52$ DH) reflects the fundamental
  difference between volume-averaged abundance (ODE) and 2-D areal coverage
  (COMSTAT), not a model failure.}
  \label{fig:hamilton_vs_comstat}
\end{figure}

\section{Towards a Unified Bayesian-Metabolic Framework}
\label{sec:integration_unified}

The two inference frameworks are, at present, methodologically disjoint: the
in-vitro analysis (Chapter~\ref{ch:heine}) produces a full Bayesian posterior but
uses only an uninformative, iteratively narrowed prior, while the in-vivo analysis
(Chapter~\ref{ch:dieckow}) injects the metabolic sign prior but collapses to a point
estimate. Their natural synthesis is a single scheme that does both --- a
\emph{metabolically-regularised TMCMC posterior}.

\paragraph{A unified objective.}
Concretely, the metabolic sign penalty $\calP(\bA;\bF)$ of
Eq.~\eqref{eq:dieckow_penalty} is a log-prior: $\log p_{\mathrm{sign}}(\bA) =
-\calP(\bA;\bF)$, the half-normal density that is flat for correct-sign entries and
quadratically suppresses sign violations in proportion to the metabolic evidence
$|F_{ij}|$. Adding it to the TMCMC prior of Section~\ref{sec:theory_tmcmc} gives the
tempered target
\begin{equation}
  p_m(\btheta) \;\propto\; p(\yobs\mid\btheta)^{\beta_m}\;
    \underbrace{p_0(\btheta)\,
    \exp\!\bigl[-\calP(\bA;\bF)\bigr]}_{\text{metabolically-informed prior}},
  \label{eq:unified}
\end{equation}
so the sampler explores the \emph{full} posterior --- credible intervals,
multimodality, model evidence --- while the metabolic prior supplies exactly the
identifiability that sparse clinical data lack. This directly targets the principal
limitation of Chapter~\ref{ch:dieckow} (the AGORA improvement was only marginally
significant at $N=10$, $p\!\approx\!0.07$): with posterior uncertainty made explicit,
the question shifts from a single point estimate to whether the metabolic prior
sharpens the posterior on the constrained entries, a quantity that remains
informative even in small cohorts. We ran exactly this coupling on the in-vitro
data --- a sign-prior-regularised TMCMC ($\sigma_{\mathrm{KEGG}}=0.15$ commensal,
$0.08$ dysbiotic) whose MAP matrix is scored against a KEGG/eHOMD net-flow matrix.
The result is informative precisely because it is \emph{asymmetric}: the metabolic
prior raises the MAP sign agreement in the commensal conditions (CS $56\!\to\!80\%$,
CH $25\!\to\!81\%$) but lowers it in the dysbiotic ones (DS $88\!\to\!44\%$,
DH $53\!\to\!44\%$). The prior aligns with health, where metabolite-mediated
mutualism organises the community, yet conflicts with dysbiosis, where the data
demand sign reversals --- concentrated on the \textit{F.~nucleatum} edges, where
shear-driven competition overrides metabolic mutualism --- that the caries-context
metabolic network does not encode. This recovers, in vitro, the very pattern the
in-vivo analysis made central: interaction sign is metabolically predictable in
health, but in dysbiosis the data must be allowed to override the prior.

\paragraph{Scaling and personalisation.}
With larger clinical series the Bayesian machinery quantifies heterogeneity a point
estimate hides; a hierarchical extension --- shared $\bA^{\mathrm{pop}}$ with
patient-specific deviations --- would partition community ecology from individual
variation. Personalising the prior via patient-specific MICOM fluxes $\bF^{(p)}$
would turn sign regularisation into a route toward individualised dysbiosis prediction.

\section{Spatial Extension: Reaction--Diffusion PDE}
\label{sec:integration_pde}

The $0$D model of Chapter~\ref{ch:heine} describes a well-mixed material point.
To resolve the depth stratification seen in confocal (FISH) imaging, we promote the
species fractions to fields $\phi_i(z,t)$ over the biofilm thickness $z\in[0,L]$ and
add molecular transport to the Hamilton reaction, giving a reaction--diffusion--advection
system in the long tradition of reaction--diffusion models of pattern formation and
biofilms~\cite{Turing1952Morphogenesis, Murray2003MathBiologyII, Kreft2001IbmBiofilm, Ward2003MathBiofilm}.

\subsection{Strong form}
\label{sec:pde_strong}
Each species field obeys
\begin{equation}
  \frac{\partial \phi_i}{\partial t}
  = \underbrace{D_i\,\frac{\partial^2 \phi_i}{\partial z^2}}_{\text{diffusion}}
    \;-\; \underbrace{u\,\frac{\partial \phi_i}{\partial z}}_{\text{advection}}
    \;+\; \underbrace{R_i(\bphi,\bpsi)}_{\text{NSP reaction}},
  \qquad i = 1,\dots,n,
  \label{eq:pde_strong}
\end{equation}
where $R_i$ is the per-point NSP reaction of Section~\ref{sec:theory_replicator}
(carrying the viability $\psi_i$, void fraction $\phi_0$, and constraint multiplier
$\gamma$), $D_i$ are species diffusivities, and $u$ is a common advective drift towards
the bulk. Diffusion and advection act on $\phi_i$ only; the reaction acts on the full NSP
state $\mathbf{g}(z)=(\phi_1,\dots,\phi_n,\phi_0,\psi_1,\dots,\psi_n,\gamma)$ at each $z$.

\paragraph{Boundary and initial conditions.}
At the substratum a no-flux (homogeneous Neumann) condition $\partial_z\phi_i|_{z=0}=0$
expresses the impermeable surface; at the biofilm--bulk interface a Dirichlet condition
$\phi_i(L,t)=\phi_{\mathrm{bulk},i}$ fixes the planktonic composition (taken as the
Day-1 median). The diffusivities $D_i$ and drift $u$ are structural parameters of
the model; with only $T=5$ FISH timepoints available (Days 1, 6, 10, 15, 21),
the temporal derivative $\partial_t\phi_i$ is too coarsely sampled to identify $D_i$
quantitatively by finite differencing~\cite{Stewart1998Diffusion}.
A PINN inverse fit to the same profiles (Section~\ref{sec:pde_fish}) yields
$D_i \approx 8\times10^{-4}$ (DH) and $D_i \approx 8\times10^{-3}$ (CH) in
normalised units, with advection $u$ reversing sign between conditions.

\subsection{Numerical scheme}
\label{sec:pde_numerics}
Equation~\eqref{eq:pde_strong} is discretised by the \emph{method of lines} with
\emph{Lie operator splitting} of the stiff reaction from the transport:
\begin{equation}
  \bphi^{\,n+1} \;=\; \mathcal{T}_{\Delta t}\bigl(\mathcal{R}_{\Delta t}(\bphi^{\,n})\bigr),
  \label{eq:splitting}
\end{equation}
a first-order ($\mathcal{O}(\Delta t)$) splitting.

\paragraph{Reaction substep $\mathcal{R}$.}
At every grid node the NSP material-point system is advanced by \emph{implicit Euler}
and solved with Newton's method ($n_{\mathrm R}=5$ iterations, internal step
$\Delta t_h=10^{-3}$), vectorised across nodes (one \texttt{vmap} over the $z$-axis).
Implicit integration is chosen because the reaction is stiff near the simplex boundary,
where the logarithmic barrier ($K_p=10^{-4}$) stiffens the Jacobian.

\paragraph{Transport substep $\mathcal{T}$ (diagonal baseline).}
Diffusion uses a second-order central stencil; advection uses first-order upwind, giving a
globally first-order ($\mathcal{O}(\Delta z)$) spatial scheme. The no-flux wall is imposed
by a ghost node; the Dirichlet bulk value is written into the last node. Explicit Euler
advances the transport; the numerical diffusion $\tfrac12 u\,\Delta z$ must be kept below
the physical $D_i$, which constrains $\Delta z$.

\subsection{Simplex preservation: volume-filling cross-diffusion}
\label{sec:pde_xdiff}
The diagonal scheme diffuses each $\phi_i$ independently and then closes the simplex by
re-setting $\phi_0=1-\sum_i\phi_i$ and clipping to $[0,1]$; this is a deliberate but
\emph{non-conservative} approximation. A simplex-preserving formulation replaces it with a
\emph{volume-filling} (size-exclusion) cross-diffusion flux~\cite{PainterHillen2002VolumeFilling}
\begin{equation}
  J_i = -D_i\Bigl[(1-\rho)\,\partial_z\phi_i + \phi_i\,\partial_z\rho\Bigr],
  \quad \rho=\sum_{j}\phi_j,
  \qquad
  \frac{\partial\phi_i}{\partial t} = -\partial_z J_i - u\,\partial_z\phi_i + R_i,
  \label{eq:xdiff}
\end{equation}
discretised conservatively by finite volumes (fluxes evaluated on cell faces). The flux is
\emph{degenerate}, $J_i\to0$ as $\rho\to1$ or $\phi_i\to0$, so $0\le\phi_i,\rho\le1$ are
preserved \emph{by construction} and $\phi_0=1-\rho$ is a genuine void fraction rather than
a patched residual. Crucially, cross-coupling enters only through the shared crowding
gradient $\partial_z\rho$, so the model retains exactly $n$ free self-mobilities $D_i$
(no $n^2$ blow-up) and remains identifiable from $z$-profiles. Both schemes share the same
sign structure and fixed points; the cross-diffusion form is used wherever conservation and
boundedness matter.

\subsection{Stability, conservation, and verification}
\label{sec:pde_verification}
The explicit transport substep is conditionally stable. The diffusive and advective
Courant--Friedrichs--Lewy (CFL) limits
\begin{equation}
  \Delta t \;\le\; \tfrac12\,\frac{\Delta z^2}{\max_i D_i}
  \qquad\text{and}\qquad
  \frac{u\,\Delta t}{\Delta z} \;\le\; 1
  \label{eq:cfl}
\end{equation}
must both hold; the cross-diffusion solver enforces the diffusive bound automatically
(clamping $\Delta t$ to $0.9$ of the limit), and production runs are reported at a fixed
sub-CFL $\Delta t$ satisfying both.

Verification in a closed-box setting (reaction off, $u=0$) confirms $\mathcal{O}(\Delta z^2)$ interior convergence, $\mathcal{O}(\Delta z)$ boundary order at the ghost-node wall, and exact mass conservation for the finite-volume cross-diffusion scheme; the diagonal scheme drifts by $+2.2\times10^{-3}$ in the sub-critical regime and overshoots $\rho>1$ when crowded. Full convergence plots and conservation table are in Appendix~\ref{app:pde}.

\paragraph{Qualitative model--data comparison.}
Figure~\ref{fig:pde_ch_profile} shows the CH attractor depth profile from the
cross-diffusion PDE alongside the corresponding FISH $z$-profile (Day~1 median,
normalised), using diffusivities fit to the commensal-HOBIC data
($D_{\mathrm{CH}}$, loss $\approx 0.015$, volume-filling cross-diffusion scheme,
$N_z=32$).  The model reproduces the qualitative features: \textit{S.~oralis}
dominates the outer layer and \textit{P.~gingivalis} is confined near the
substrate, consistent with the depth-ordering and Manders results of
Section~\ref{sec:pde_fish}.  Quantitative agreement is limited by the coarse
temporal sampling ($T=5$ FISH timepoints), which prevents full identification of
the $D_i$ (Section~\ref{sec:pde_strong}); fitted values should be read as
order-of-magnitude estimates rather than measured transport coefficients.

\begin{figure}[htbp]
  \centering
  \includegraphics[width=0.74\textwidth]{fig_pde_ch_profile.pdf}
  \caption[PDE depth profile vs.\ FISH observation (CH attractor)]{%
  Qualitative comparison of the reaction--diffusion model and the FISH
  $z$-profile for the commensal-HOBIC (CH) attractor.
  \textbf{(a)}~FISH median depth profile (Day~1, $N_z=32$ bins, normalised
  relative abundance). \textbf{(b)}~Steady-state $z$-profile from the
  volume-filling cross-diffusion PDE with diffusivities fit to the CH
  $z$-profiles (loss $= 0.015$).
  The outer-layer dominance of \textit{S.~oralis} (blue) and inner
  sequestration of \textit{P.~gingivalis} (red) are reproduced qualitatively;
  quantitative fits await continuous live-imaging data for full $D_i$
  identification.}
  \label{fig:pde_ch_profile}
\end{figure}

\subsection{Empirical depth structure from 3-D FISH imaging}
\label{sec:pde_fish}
The spatial model is ultimately answerable to the real depth structure of the
biofilm. We extracted the full three-dimensional $(x,y,z)$ organisation of the five
species from the Heine HOBIC FISH dataset --- $84$ titanium fields of view across
all $11$ confocal \texttt{.lif} acquisitions, decoded \emph{per voxel} (without the
usual $xy$-averaging) into the five-species channels. Two findings, both resting on
billions of voxels, emerge and are the principal empirical result of the spatial
work; the reaction--diffusion model is the methodology that confronts them.
As a sanity check, Figure~\ref{fig:fish_vs_experiment} confirms that the
$z$-integrated FISH composition at day~21 agrees with the Heine bulk culture
measurements, validating the voxel decoding pipeline.

\begin{figure}[htbp]
  \centering
  \includegraphics[width=0.68\textwidth]{fig_fish_vs_experiment.pdf}
  \caption[FISH voxel composition vs Heine culture experiment]{%
  Species fractions from $z$-integrated FISH voxel counts (day~21, all FOVs)
  versus the Heine et al.\ culture-experiment replicate medians (day~21 HOBIC).
  Good agreement validates the voxel-level channel decoding; residuals reflect
  the difference between single-FOV local composition and the bulk-culture mean.}
  \label{fig:fish_vs_experiment}
\end{figure}

\begin{figure}[htbp]
  \centering
  \includegraphics[width=0.78\linewidth]{fig_fish_clsm_feature.png}
  \caption[HOBIC FISH: top-down view and depth cross-section]{%
  Single-FOV CLSM summary (commensal-HOBIC, Day~1).
  \textbf{Left:} $xy$ maximum-intensity projection (MIP); the five-species composite
  shows a patchy, microcolonial architecture dominated by
  \textit{S.~oralis} (blue) and \textit{Vd/Vp} (yellow).
  \textbf{Right:} $xz$ depth cross-section MIP; the vertical extent of
  each species reveals the stratification structure that motivates the
  reaction--diffusion model of Section~\ref{sec:integration_pde}.}
  \label{fig:fish_clsm_feature}
\end{figure}

\paragraph{\textit{Fn}--\textit{Pg} decoupling is a depth separation.} In the
dysbiotic-HOBIC condition the lateral ($xy$-projected) Manders overlap
coefficient~\cite{Manders1993Colocalization} of
\textit{F.~nucleatum} and \textit{P.~gingivalis} is near unity --- the two species
occupy the same lateral footprint --- yet their true \emph{three-dimensional}
voxel-wise overlap collapses to $M_1\approx0.22$--$0.33$ from Day~6 onward (versus
$0.4$--$0.7$ in the commensal-HOBIC control; Figure~\ref{fig:fish_coloc}). The two
species therefore share an $(x,y)$ position but sit at \emph{different depths}: \textit{Pg}
settles \emph{beneath} \textit{Fn}. This is a direct 3-D confirmation of \textit{Pg}
deep-layer sequestration, invisible to any 2-D (depth-collapsed) analysis.

\paragraph{Fn--Pg co-aggregation collapses by Day~6 in DH.}
To quantify \emph{when} the depth separation emerges, we compute the pair
cross-correlation function $g(r)$ between \textit{F.~nucleatum} (Fn) and
\textit{P.~gingivalis} (Pg) voxels following Valm et al.~\cite{Valm2012PairCorrelation}:
$g(r) > 1$ indicates co-aggregation at distance~$r$, $g(r) = 1$ the Poisson
(random) baseline, and $g(r) < 1$ mutual avoidance.
In the dysbiotic-HOBIC condition (Figure~\ref{fig:fish_pair_correlation}b), Day~1
shows a pronounced short-range peak ($g \approx 2.8$ at $r < 2\,\mu\text{m}$):
Fn and Pg are initially co-localised, as expected immediately after inoculation.
By Day~6 the same short-range signal collapses to $g \approx 0.3$, a strong
avoidance signature, and remains below unity through Day~21. This temporal
transition marks the onset of the depth-stratification event: Pg migrates to the
anaerobic inner layer beneath Fn within the first five experimental days, and the
two species thereafter occupy spatially exclusive zones even at sub-5\,$\mu$m
scales. In the commensal-HOBIC condition (Figure~\ref{fig:fish_pair_correlation}a),
$g(r) \approx 1$ throughout, with modest co-aggregation fluctuations consistent
with the patchy, microcolonial architecture. The $g(r)$ analysis thus provides a
statistical timestamp for the dysbiotic spatial reorganisation that complements the
depth-profile and Manders data.

\begin{figure}[htbp]
  \centering
  \includegraphics[width=0.82\linewidth]{fish_pair_correlation.pdf}
  \caption[Fn--Pg pair cross-correlation $g(r)$]{Pair cross-correlation function
  $g(r)$ between \textit{F.~nucleatum} and \textit{P.~gingivalis} voxels across
  days for commensal-HOBIC (\textbf{a}) and dysbiotic-HOBIC (\textbf{b}).
  In DH, $g$ collapses from $>2.5$ (Day~1, co-aggregation) to $<0.5$ (Day~6,
  avoidance) at short range ($r < 5\,\mu$m), marking the onset of depth
  stratification. CH shows $g \approx 1$ throughout. Computed from 3-D FISH voxel
  stacks ($\phi_i > 0.15$ threshold, $8\times10^4$ random pairs per bin;
  Valm et al.\ 2012).}
  \label{fig:fish_pair_correlation}
\end{figure}

\paragraph{Dysbiosis is laterally homogeneous.} Quantifying lateral patchiness by
the per-voxel coefficient of variation (CV) of each species across the $xy$ plane,
the dysbiotic-HOBIC biofilm is markedly \emph{more homogeneous} than the commensal
one (Figure~\ref{fig:fish_lateral}): the across-species mean lateral CV is
$0.75$ (commensal, $n=67$ FOV) versus $0.52$ (dysbiotic, $n=17$ FOV), a difference
significant for every species (Mann--Whitney $\mathrm{CH}\!>\!\mathrm{DH}$,
$p=2.2\times10^{-4}$ pooled; per species \textit{So}/\textit{An}/\textit{Fn}
$p<0.001$, \textit{Veillonella} $p=0.003$, \textit{Pg} $p=0.041$). \textit{P.\
gingivalis} shows the weakest homogenisation, retaining the most patchiness. The
commensal biofilm is, by contrast, organised into discrete microcolonies
(CV $0.4$--$0.9$). Taken together: \emph{dysbiosis = a laterally confluent,
vertically stratified mat with \textit{Pg} sequestered at depth; commensal health =
a patchy, microcolonial structure}.

\begin{figure}[htbp]
  \centering
  \includegraphics[width=0.72\textwidth]{fig_fish_fnpg_coloc_thesis.pdf}
  \caption[3-D \textit{Fn}--\textit{Pg} colocalisation]{\textit{Fn}--\textit{Pg}
  colocalisation from 3-D FISH. The lateral ($xy$-projected) Manders overlap is near
  unity, but the true voxel-wise 3-D overlap drops to $M_1\approx0.22$--$0.33$ in
  dysbiotic-HOBIC from Day~6 --- the two species share a lateral footprint but
  separate in depth (\textit{Pg} below \textit{Fn}).}
  \label{fig:fish_coloc}
\end{figure}

\begin{figure}[htbp]
  \centering
  \includegraphics[width=0.72\textwidth]{fig_fish_lateral_heterogeneity_thesis.pdf}
  \caption[Lateral homogenisation in dysbiosis]{Lateral heterogeneity (per-voxel
  $xy$ coefficient of variation) by species and condition. Dysbiotic-HOBIC is
  significantly more homogeneous than commensal-HOBIC for every species (pooled
  Mann--Whitney $p=2.2\times10^{-4}$, $84$ FOV); \textit{Pg} homogenises least.}
  \label{fig:fish_lateral}
\end{figure}

\begin{figure}[htbp]
  \centering
  \includegraphics[width=0.72\linewidth,height=0.50\textheight,keepaspectratio]{fig_fish_zprofile_temporal.pdf}
  \caption[Temporal evolution of biofilm depth profiles]{Biofilm species composition
  as a function of depth ($z=0$: substrate) across days 1--21 for commensal-HOBIC (CH,
  top) and dysbiotic-HOBIC (DH, bottom).  Stratification intensifies over time;
  by day~6 \textit{S.~oralis} (blue) occupies the outer layers while
  \textit{P.~gingivalis} (red) retreats toward the substrate in CH, and the DH
  profiles become progressively more species-uniform at each depth --- consistent
  with the lateral homogenisation result (Figure~\ref{fig:fish_lateral}).}
  \label{fig:fish_zprofile_temporal}
\end{figure}

\begin{figure}[htbp]
  \centering
  \includegraphics[width=0.82\linewidth]{fish_cooccurrence_depth.pdf}
  \caption[Spatial co-occurrence vs.\ predicted interaction sign]{Spatial co-occurrence
  analysis: Pearson $r$ between species-pair voxel intensities at each normalised
  depth, with Pg-involved pairs highlighted. \textbf{(Left)} CH and DH depth profiles
  for So/An/Vd-Vp/Fn vs.\ Pg; negative $r$ (anti-correlation) builds in the Pg zone
  (shaded) consistent with competitive exclusion. \textbf{(Centre)} Pg-involved pairs
  agree with the ODE interaction sign in 4/4 CH pairs ($p=0.06$, binomial) vs.\ 2/4 DH
  pairs; non-Pg pairs agree at chance. \textbf{(Right)} Day-by-day agreement: signal
  strengthens as the biofilm matures.}
  \label{fig:fish_cooccurrence}
\end{figure}

\paragraph{\textit{S.~oralis}--\textit{Veillonella} depth ordering confirms metabolic cross-feeding.}
The AGORA metabolic backbone predicts lactate flux from \textit{S.~oralis} (producer)
to \textit{Veillonella} spp.\ (consumer), implying the consumer should sit closer to
the bulk to intercept outward-diffusing substrate.
We computed species centre-of-mass depth $\bar{z}_i$ for each condition--day
combination ($n=10$): \textit{Veillonella} is shallower than \textit{S.~oralis} in
all 10 timepoints ($\Delta z = -2.81\,\mu\mathrm{m}$, $p=0.002$,
Figure~\ref{fig:welch_so_vei_depth}), consistent with Welch et al.\
\cite{Welch2016Biogeography} and constituting a geometry-based validation of the
So$\to$Vei cross-feeding edge.

\begin{figure}[htbp]
  \centering
  \includegraphics[width=0.74\textwidth, height=0.30\textheight, keepaspectratio]{fig_welch_so_vei_depth.pdf}
  \caption[Metabolic depth ordering: \textit{S.~oralis} deeper than \textit{Veillonella}]{%
  Spatial cross-feeding signature: \textit{Veillonella} (lactate consumer) sits
  consistently shallower than \textit{S.~oralis} (lactate producer).
  \textbf{(a)}~Species centre-of-mass depth across days for CH (circles, solid)
  and DH (squares, dashed); shading marks the consumer--producer gap.
  \textbf{(b)}~Signed depth difference $\Delta z = z_\mathrm{Vei} - z_\mathrm{So}$
  per timepoint; negative in 10/10 cases (two-sided sign test $p=0.002$, $n=10$).
  \textbf{(c)}~Comparison with Welch et al.\ (2016), who reported interdigitated
  \textit{Veillonella}--\textit{Streptococcus} co-localisation in supragingival
  CLASI-FISH; our signed depth metric ($-2.81\,\mu$m mean) provides a quantitative
  extension to subgingival titanium biofilm.}
  \label{fig:welch_so_vei_depth}
\end{figure}

\paragraph{All five species: \textit{Veillonella} is the community's shallowest member.}
Extending the CoM analysis to all five species (Figure~\ref{fig:depth_ordering_all5sp}),
\textit{Veillonella} is the shallowest in every timepoint,
significantly above So ($p=0.001$), An ($p=0.006$), and Fn ($p=0.006$);
the remaining four species cluster at $25$--$27\,\mu\mathrm{m}$ without significant
pairwise ordering.
The So$\to$Vei and An$\to$Vei edges predicted by AGORA are both confirmed; the
Fn$\to$Vei separation is an additional finding not encoded in the FBA network.

\begin{figure}[htbp]
  \centering
  \includegraphics[width=0.85\linewidth, height=0.29\textheight, keepaspectratio]{fig_depth_ordering_all5sp.pdf}
  \caption[All-species centre-of-mass depth ordering]{Centre-of-mass depth ranking
  for all five species ($n=10$ condition--day combinations).
  \textbf{(a)}~Mean $\pm$ s.d.\ CoM depth per species (sorted shallow to deep).
  \textbf{(b)}~Pairwise depth difference matrix
  $\Delta z = \bar{z}_\mathrm{col} - \bar{z}_\mathrm{row}$; stars mark
  Wilcoxon-significant pairs (* $p<0.05$, ** $p<0.01$, *** $p<0.001$).
  \textbf{(c)}~Schematic depth cartoon with arrows indicating significant pairs.
  \textit{Veillonella} is the shallowest species in all three representations,
  separated from all others at $p \le 0.006$.}
  \label{fig:depth_ordering_all5sp}
\end{figure}

\paragraph{PINN-inferred transport parameters.}
To move beyond order-of-magnitude placeholders, a physics-informed neural network
(PINN) was fitted to the FISH $\phi_i(z,t)$ profiles using the same PDE operator
(Eq.~\eqref{eq:pde_strong}), with $D_i$ and $u$ as jointly trainable parameters
alongside the network weights \cite{Karniadakis2021PINN}.
The inferred diffusivities are $D_i \approx 8\times10^{-4}$ (normalised) in DH and
$D_i \approx 8\times10^{-3}$ in CH --- a tenfold reduction in matrix permeability
under dysbiotic conditions that is uniform across all five species
(Figure~\ref{fig:pinn_diffusivity}).
The dominant condition contrast lies instead in the advection sign:
$u_\mathrm{DH} = -1.76$ (inward, toward the substratum) versus
$u_\mathrm{CH} = +3.02$ (outward, toward the bulk).
Depth sequestration of \textit{P.~gingivalis} in dysbiosis is therefore
advection-driven rather than reflecting a species-specific diffusivity deficit;
globally denser matrix and reversed bulk flow jointly account for the stratification
observed in FISH.

\begin{figure}[htbp]
  \centering
  \includegraphics[width=0.78\textwidth]{fig_pinn_diffusivity.pdf}
  \caption[PINN-inferred diffusivities and advection: DH vs CH]{%
  Physics-informed neural network (PINN) transport parameter estimates.
  \textbf{(a)}~Observed FISH depth profiles for Dysbiotic HOBIC (day 21),
  showing the late-stage Pg enrichment at depth.
  \textbf{(b)}~PINN-inferred diffusivities $D_i$ for all five species under DH
  (red) and CH (blue) conditions; all species show a ${\sim}10$-fold reduction
  in DH. The advection parameter $u$ reverses sign between conditions
  ($u_\mathrm{DH}=-1.76$, $u_\mathrm{CH}=+3.02$, normalised units),
  indicating inward bulk flow in dysbiosis as the primary driver of
  \textit{Pg} depth sequestration.}
  \label{fig:pinn_diffusivity}
\end{figure}

\paragraph{Interpretation of transport parameters.}
The tenfold global reduction in $D_i$ in DH is consistent with the denser,
more structured extracellular matrix of a mature dysbiotic biofilm: higher
exopolysaccharide content and closer species packing physically restrict molecular
diffusion for all guild members equally, without invoking any species-specific
mechanism.
The advection term $u$ captures net bulk fluid flow relative to the growing biofilm
front; its sign reversal ($u<0$ inward in DH, $u>0$ outward in CH) may reflect
the anaerobic pocket dynamics of the DH condition, where oxygen depletion at depth
creates a nutrient gradient that drives net inward transport of solutes and,
secondarily, the microbial community it supports.
The finding that $D_\mathrm{Pg}$ is not distinctively low compared with the other four
species implies that spatial sequestration of \textit{P.~gingivalis} at the
substratum is not a passive diffusion barrier effect but an active advective process,
consistent with the known motility-independence of \textit{Pg} and its reliance on
community structure for niche establishment \cite{Hajishengallis2014Polymicrobial}.

\paragraph{What the model captures --- and what it does not.} Confronted with these
data, the reaction--diffusion model of this section reproduces the dysbiotic lateral
homogenisation: under reaction-dominated dynamics with a laterally uniform boundary,
any initial lateral structure decays (the cross-diffusion variant retains only
marginally more). It does \emph{not} reproduce the commensal microcolonial
patchiness --- a continuum reaction--diffusion field has no mechanism to sustain
discrete colonies against diffusive smoothing. The honest reading is that the
\emph{depth} stratification and the dysbiotic homogenisation are within the model's
reach, whereas commensal patchiness points beyond it, to stochastic colonisation or
adhesion-driven pattern formation not encoded in Eq.~\eqref{eq:pde_strong}. The
scientific weight here sits on the data; the PDE is the lens that makes precise
which features it can and cannot explain.

\subsection{Scope and limits of the mean-field description}
\label{sec:discussion_mismatch}

The ODE/COMSTAT mismatch noted above --- competitive exclusion in the volume-averaged
model against near-uniform 2-D coverage, reconciled by the vertical niche partitioning of
Section~\ref{sec:pde_fish} --- marks the boundary of the mean-field model's applicability and
is the primary motivation for the spatial PDE extension of this chapter. Full closure would
require jointly inferring $\mathbf{A}$ from temporal abundance data and spatial coverage
metrics simultaneously.

\section{Mechanical Stratification: Finite-Element Realisation}
\label{sec:integration_fem}

The spatial analysis of Section~\ref{sec:integration_pde} establishes \emph{where}
each species resides --- the depth-resolved composition $\phi_i(z)$ and, in
particular, the substratum sequestration of \textit{P.~gingivalis}. It leaves open
the mechanical question that the Extended Hamilton Principle of
Chapter~\ref{ch:theory} was, in its biofilm-mechanics reading, built to answer: what
\emph{stress state} does a given community structure imprint on the film? Because
the local biomass fraction sets the local growth, a depth-varying composition drives
a depth-varying growth and hence --- once the film is constrained by the implant
surface --- a depth-varying residual stress. This section realises that chain
numerically, embedding the composition field as the growth driver of a finite-element
continuum model and reading off the mechanical stratification it produces. It is the
mechanical completion of the ``Spatial Stratification'' the title promises:
composition depth-structure $\rightarrow$ mechanical state $\rightarrow$ clinic.

\subsection{Coupling concept and constitutive model}
\label{sec:fem_method}

The biofilm material model is invoked \emph{at each Gauss point}, replacing the
phenomenological constitutive law of a host finite-element framework with the
composition-driven growth mechanics derived here. Kinematically we adopt the
classical multiplicative growth split~\cite{Rodriguez1994Growth}
\begin{equation}
  \bF = \bF_{\!e}\,\bF_{\!g},
  \qquad
  \bF_{\!g} = \operatorname{diag}\!\bigl(1,\,1,\,\lambda_z(\phi_\mathrm{Pg})\bigr),
  \label{eq:growth_split}
\end{equation}
in which the inelastic factor $\bF_{\!g}$ describes stress-free volumetric growth and
the elastic factor $\bF_{\!e}=\bF\bF_{\!g}^{-1}$ carries the stress. Its magnitude is
set point-wise by the local \textit{P.~gingivalis} fraction $\phi_\mathrm{Pg}(z)$ from
the ecological model; its kinematic form follows the constraint. For \emph{unconfined}
thickening --- the free-swelling calibration of Section~\ref{sec:fem_growth} --- growth
is anisotropic and substratum-normal, $\bF_{\!g}=\operatorname{diag}(1,1,\lambda_z)$ as
in Eq.~\eqref{eq:growth_split}, new biomass accreting away from the implant surface. For
the laterally \emph{confined} residual-stress column of Section~\ref{sec:fem_residual},
however, a purely substratum-normal growth would relax freely against the traction-free
top and develop \emph{no} in-plane stress; the in-plane residual stress is generated by
the confinement of \emph{lateral} growth, so there growth is modelled as isotropic,
$\bF_{\!g}=J_g^{1/3}\bI$, with the volume factor $J_g$ matched to the same CLSM volume
increase $J_g=\lambda_z(\phi_\mathrm{Pg})$. Because the true growth anisotropy is not
measured, this volume-matched isotropy is an explicit modelling choice; it fixes the
residual-stress \emph{shape} from $\phi_\mathrm{Pg}(z)$ but leaves its absolute
\emph{scale} model-dependent (Section~\ref{sec:fem_validation}). The elastic response
uses a compressible neo-Hookean potential
\cite{Holzapfel2000}, giving the Cauchy stress
\begin{equation}
  \bsigma
  = \frac{\mu}{J_e}\bigl(\bb_e-\bI\bigr)
  + \frac{\lambda}{J_e}\,\ln J_e\,\bI,
  \qquad
  \bb_e=\bF_{\!e}\bF_{\!e}^{\!\top},\;\;J_e=\det\bF_{\!e},
  \label{eq:neo_hookean}
\end{equation}
with the consistent material tangent supplied analytically to preserve the quadratic
convergence of the global Newton solve. The neo-Hookean elastic factor is a
deliberate \emph{fast-elastic} idealisation: over the seconds-to-minutes timescale on
which the depth structure is mechanically interrogated, biofilms respond essentially
elastically~\cite{Stoodley2002Material}, so Eq.~\eqref{eq:neo_hookean} returns the
short-time stress \emph{before} viscous relaxation (the long-time complement is
treated in Section~\ref{sec:fem_validation}).

Two coupling implementations were realised. A self-contained user-material
subroutine communicates with the model over a local socket (a Fortran
\texttt{ISO\_C\_BINDING} client, requiring no external link step), which is
convenient for prototyping; for production and for portability to a second solver the
growth driver $\phi_\mathrm{Pg}(z)$ is instead delivered as a \emph{predefined field
variable}, evaluated per Gauss point from the element coordinates. The latter path
contains no inter-process communication and maps one-to-one onto the field-variable
mechanism of other commercial solvers, so the constitutive routine is portable
essentially unchanged.

\paragraph{Material parameters, boundary conditions, and loading.}
The coupled implant--tooth--bone models of Section~\ref{sec:fem_tooth} use the
linear-elastic isotropic parameter set of Table~\ref{tab:fem_materials}, in
consistent mm--N--MPa units (stresses in MPa). The values are standard for dental
implant finite-element analysis \cite{Geng2001,Sevimay2005}; the PDL modulus is the
linear secant value of \citet{ReesJacobsen1997} and the dentin modulus follows
\citet{Kinney2003}. The periodontal ligament (PDL) and the biofilm are
deliberate linear idealisations of nonlinear/uncertain materials (the true PDL is
nonlinear/hyperelastic \cite{Cattaneo2005}), so every
biofilm-driven result is reported as a dimensionless dysbiosis \emph{ratio}
(DH/CH$\,\approx\,2.3$--$3.3\times$ across all models) rather than an absolute stress,
and the PDL enters only through the implant-versus-tooth \emph{contrast}. The crown
stiffness is near-rigid relative to the load path: it scarcely affects the
peri-implant bone stress, which is governed by the load resultant and its moment arm.
The bone is held fully fixed on the artificial crop faces ($U_1{=}U_2{=}U_3{=}0$; a
St.-Venant truncation, verified non-artefactual by a periodic representative-volume
closure). Osseointegration is modelled as a tied interface (\texttt{*TIE},
\texttt{ADJUST=NO}), the crown is tied to the abutment top, and a frictional general
contact ($\mu=0.3$) is used only for the de-integrated micromotion variant. Biofilm
growth is imposed as an isotropic volumetric eigenstrain ($\varepsilon_\mathrm{g}=0.19$,
thermal analogy) in a first step; the occlusal bite (a $100\,$N resultant at
$30^\circ$ to the axis, per ISO~14801) is applied in a second step, on the abutment
top for the bare case and --- crucially --- on the \emph{crown occlusal table} once a
restoration is present, which lengthens the moment arm (Section~\ref{sec:fem_tooth}).
The large coupled assembly uses linear tetrahedra (C3D4, ${\sim}2.7\times10^{5}$
elements); the ISO~14801 design and bone-loss coupons use quadratic tetrahedra (C3D10),
since linear elements under-predict the thread-root concentration by $9$--$15\%$. Mesh
convergence on the coupon thread-root peak is within $\pm7\%$.

\begin{table}[htbp]
  \centering\small
  \begin{tabular}{@{}lrrl@{}}
    \toprule
    \textbf{Material} & $E$ & $\nu$ & \textbf{Role / basis}\\
     & (MPa) & & \\
    \midrule
    Titanium (Ti-6Al-4V)        & $110\,000$ & $0.34$ & implant fixture + abutment\\
    Cortical bone               & $13\,700$  & $0.30$ & dense shell / lamina dura ($1.8$\,mm)\\
    Cancellous bone             & $1\,000$   & $0.30$ & trabecular core (type III--IV)\\
    Dentin                      & $18\,000$  & $0.31$ & natural neighbour tooth\\
    PDL                         & $50$       & $0.45$ & $0.25$\,mm ligament (linear idealisation)\\
    Crown (lithium-disilicate)  & $95\,000$  & $0.30$ & load-bearing restoration ($210\,000$ for zirconia)\\
    Gingiva (mucosa)            & $3$        & $0.45$ & peri-implant soft-tissue cuff\\
    Biofilm                     & $1.0$      & $0.45$ & growth collar ($+\,\varepsilon_\mathrm{g}=0.19$; modulus uncertain $\Rightarrow$ ratios)\\
    \bottomrule
  \end{tabular}
  \caption[Finite-element material parameters]{Linear-elastic isotropic material
  parameters for the coupled implant--tooth--bone models (mm--N--MPa units). PDL and
  biofilm are explicit idealisations of nonlinear/uncertain materials; biofilm-driven
  outputs are reported as dysbiosis ratios, not absolute stresses.}
  \label{tab:fem_materials}
\end{table}

\subsection{Growth kinematics from confocal data}
\label{sec:fem_growth}

The single free magnitude of Eq.~\eqref{eq:growth_split}, the growth stretch
$\lambda_z$, is calibrated against the CLSM depth data of
Section~\ref{sec:pde_fish}. Regressing the local biofilm thickening on the local
\textit{P.~gingivalis} fraction yields a clear monotone relation in the dysbiotic
condition (Figure~\ref{fig:fem_calibration}; Pearson $r=0.88$, slope
$\beta_\mathrm{depth}\approx17$), whereas the commensal control shows no comparable
thickening. The growth law is therefore data-anchored in shape: a
\textit{Pg}-enriched depth grows more, exactly where FISH places the organism. As
discussed in Section~\ref{sec:fem_validation} we treat the calibration as fixing the
\emph{shape} of $\lambda_z(\phi_\mathrm{Pg})$ while leaving its absolute scale a
single parameter $\gamma$.

\begin{figure}[htbp]
  \centering
  \includegraphics[width=0.56\textwidth]{fem_growth_calibration.pdf}
  \caption[Growth calibration from CLSM depth data]{Depth-resolved growth strain
  versus local \textit{P.~gingivalis} fraction (dysbiotic-HOBIC). The growth that
  Eq.~\eqref{eq:growth_split} requires is set by where \textit{Pg} sits: the
  thickening correlates with $\phi_\mathrm{Pg}$ at $r=0.88$
  ($\beta_\mathrm{depth}\approx17$); the commensal control shows no such trend.}
  \label{fig:fem_calibration}
\end{figure}

\subsection{Verification}
\label{sec:fem_verification}

Before any biofilm-specific result is reported, the coupled implementation is
verified against analytic expectations by four independent checks
(Table~\ref{tab:fem_verification}), in the spirit of code verification by
manufactured solutions~\cite{Roache2002MMS}. (i)~A \emph{patch test} under spatially
uniform growth returns a spatially uniform interior stress: all interior elements of
a multi-element block reproduce an \emph{identical} in-plane stress to machine
precision, the free surface carrying the expected single-element boundary layer.
(ii)~An \emph{energy-conservation} test, with growth switched off and a prescribed
stretch applied, reproduces the closed-form neo-Hookean stored energy: the external
work equals the analytic strain energy to a relative difference of
$1.3\times10^{-4}$, confirming the absence of spurious numerical dissipation and a
correct (symmetric Poisson) response. (iii)~A \emph{method-of-manufactured-solutions}
study, in which a non-trivial displacement field is forced by the corresponding
analytic body load, recovers the theoretical convergence order of the trilinear
hexahedral element: the $L^2$ displacement error decays as $h^{1.95}\approx h^2$
under uniform refinement. (iv)~The analytically supplied consistent tangent is
confirmed against a finite-difference (Kirchhoff/Jaumann) tangent, which it reproduces
to a relative error of $5\times10^{-8}$ across the full calibrated growth range,
guaranteeing the quadratic convergence of the global Newton solve. Finally, the growth
column used for the production results below is mesh-converged
(Figure~\ref{fig:fem_convergence}): under spatially uniform growth the confined
plateau stress is mesh-independent to machine precision, and for the depth-graded
column the confined interface stress at $N=16$ elements lies within $0.27\%$ of the
$N=32$ value.

\begin{table}[htbp]
  \centering\small
  \begin{tabular}{@{}llp{5.6cm}@{}}
    \toprule
    \textbf{Check} & \textbf{Quantity} & \textbf{Result}\\
    \midrule
    Patch test            & interior stress uniformity & identical interior stress (machine precision)\\
    Energy conservation   & $\lvert W_\mathrm{ext}-U_\mathrm{NH}\rvert/W_\mathrm{ext}$ & $1.3\times10^{-4}$ (no dissipation)\\
    Manufactured solution & $L^2$ order $p$            & $p=1.95\approx2$ (trilinear hexahedron)\\
    Consistent tangent    & analytic vs FD (rel.\ err.) & $5\times10^{-8}$ (full growth range)\\
    Mesh convergence      & interface stress, $N{=}16$ vs $32$ & within $0.27\%$ (uniform: exact)\\
    \bottomrule
  \end{tabular}
  \caption[Verification of the finite-element coupling]{Verification of the coupled
  growth--mechanics implementation against analytic expectations. The three code
  checks are independent of any biofilm parameter.}
  \label{tab:fem_verification}
\end{table}

\begin{figure}[htbp]
  \centering
  \includegraphics[width=0.56\textwidth]{fem_mesh_convergence.pdf}
  \caption[Mesh convergence of the growth column]{Confined in-plane stress versus the
  number of elements through the depth. Under spatially uniform growth the plateau
  stress is mesh-independent to machine precision (lower curve); for the depth-graded
  column the confined interface stress converges by $N=16$ (within $0.27\%$ of $N=32$);
  production runs use $N=16$.}
  \label{fig:fem_convergence}
\end{figure}

\subsection{Result: composition-resolved residual stress}
\label{sec:fem_residual}

The central mechanical result follows from confining the growth column to the
substrate (bonded base, traction-free top, laterally constrained), so that the lateral
growth cannot relax by free expansion and instead develops an in-plane residual stress.
Figure~\ref{fig:fem_residual} shows the resulting profile $\sigma_{xx}(z)$ together with
its compositional driver. The residual stress is \emph{not} distributed through the
depth but concentrates in a boundary layer at the substratum --- the implant interface
--- of order one column width deep: there the confined growth is held, whereas higher in
the film the traction-free surface relieves the growth by free upward expansion and the
in-plane stress decays to zero (above the layer the column simply grows taller at
$\sigma\approx0$, regardless of $\phi_\mathrm{Pg}$). \emph{Within} the interface layer
the stress tracks $\phi_\mathrm{Pg}$ monotonically --- more \textit{P.~gingivalis} drives
more confined growth and hence more compression. The condition contrast is therefore
concentrated exactly where it is clinically decisive: the dysbiotic film,
\textit{Pg}-enriched at the base, carries a far larger interface stress than the
commensal one --- at identical material parameters its peak exceeds the commensal by
$6.4\times$ and its mean interface stress by $\approx10\times$, while the commensal
profile is weak and nearly featureless. The mechanical state is thus the
\emph{composition-resolved} substratum stress that a particular dysbiotic community
imprints on the implant interface, a stress the laterally homogeneous commensal biofilm
does not develop. Because the elastic modulus of oral biofilm is uncertain by orders
of magnitude (Section~\ref{sec:fem_validation}), the stress is reported relative to
the shear modulus $\mu$, leaving the result modulus-agnostic.

\begin{figure}[htbp]
  \centering
  \includegraphics[width=0.72\textwidth]{fem_residual_stress.pdf}
  \caption[Composition-resolved substratum-interface residual stress: dysbiotic vs
  commensal]{\textbf{(left)} In-plane residual stress $\sigma_{xx}(z)/\mu$ from the
  substrate-bonded growth column and \textbf{(right)} its compositional driver
  $\phi_\mathrm{Pg}(z)$, for the dysbiotic (DH) and commensal (CH) conditions
  ($z/H=0$ at the implant interface). The stress concentrates in the substratum
  boundary layer (shaded); above it the traction-free surface relieves the growth and
  the in-plane stress decays to zero. The dysbiotic interface stress far exceeds the
  commensal (peak $6.4\times$, mean $\approx10\times$); \emph{within} the layer it
  tracks $\phi_\mathrm{Pg}$. This is the mechanical signature of the depth composition
  documented in Section~\ref{sec:pde_fish}.}
  \label{fig:fem_residual}
\end{figure}

\subsection{Clinical reading: residual stress drives delamination}
\label{sec:fem_delamination}

The residual stress acquires clinical meaning at the implant interface. Replacing the
perfectly bonded base of the growth column with a \emph{cohesive} interface (a
traction--separation law with a quadratic-stress damage-initiation criterion and
energy-based evolution) lets the growth-induced stress do mechanical work against the
film--substrate bond. The growth-induced shear concentrates at the free edge, where
cohesive damage initiates and a delamination front then propagates inward
(Figure~\ref{fig:fem_delamination}). The condition contrast is decisive: at an
\emph{identical} interface strength, the dysbiotic film delaminates over $83\%$ of
the interface against only $25\%$ for the commensal film --- a $3.3\times$ difference
whose sole cause is the \textit{Pg} depth profile that sets the growth. This is a
direct finite-element test of the dysbiosis$\rightarrow$detachment hypothesis: the
same compositional shift that the ecological model identifies as dysbiotic produces,
through its mechanical realisation, a markedly larger tendency for the biofilm to
detach from the implant --- the mechanical event that seeds peri-implant soft-tissue
exposure. It connects the spatial composition result directly to the clinical
translation of Section~\ref{sec:integration_clinical}.

\begin{figure}[htbp]
  \centering
  \includegraphics[width=0.72\textwidth]{fem_delamination.pdf}
  \caption[Growth-driven delamination: dysbiotic vs commensal]{\textbf{(left)}
  Cohesive interface damage along the film for the two conditions at full growth: the
  damage front reaches far deeper into the dysbiotic interface. \textbf{(right)}
  Delaminated fraction of the interface versus interface strength; at every strength
  the dysbiotic film detaches further, by up to $3.3\times$. Identical material and
  interface parameters --- only the $\phi_\mathrm{Pg}$ depth profile differs.}
  \label{fig:fem_delamination}
\end{figure}

\paragraph{Re-solving on an actual implant thread.} The growth column and cohesive test keep
the substratum flat, whereas the peri-implant biofilm of the Dieckow cohort grows on a
threaded, micro-rough titanium abutment. Re-solving the same CLSM-calibrated growth on a smooth
threaded titanium cross-section --- a base-bonded biofilm patch in plane strain, the
volume-matched growth imposed as a thermal eigenstrain so that no user material is needed, with a
mesh-convergent interface stress (the interior peak, with the bonded-base/free-edge corner
singularity excluded, bounded within $\pm7\%$ for $\ge24$ elements per thread period;
Figure~\ref{fig:fem_implant}B) --- makes the implant thread the central geometric feature of the
result. First, the thread is what \emph{generates} the interface stress --- the driver of thin-film
delamination from the substrate~\cite{HutchinsonSuo1992}: a flat titanium surface
lets the laterally-relaxed growth expand almost stress-free, whereas the thread frustrates it,
and the interior interface stress rises monotonically with thread depth --- roughly an order of
magnitude from flat to a deep thread (Figure~\ref{fig:fem_implant}C) --- concentrating at the
bonded thread roots (Figure~\ref{fig:fem_implant}A). Second, the dysbiosis contrast rides on top
of this geometric stress and is preserved: the dysbiotic interior interface von~Mises stress is
$\approx\!2.5\times$ the commensal, the same ordering as the flat-column residual-stress and
cohesive-delamination contrasts (Sections~\ref{sec:fem_residual}--\ref{sec:fem_delamination}).
This is not a finite-size artefact of the free patch edges: a periodic-RVE unit cell --- a single
thread period with left--right periodicity, the laterally-infinite-film limit, which removes the
free edge entirely so the whole bonded interface is interior --- gives the same picture, the flat
interface remaining essentially stress-free while the thread generates the interface stress
($\mathrm{DH/CH}=2.3\times$ in the RVE). The thread is therefore the geometric origin of the
interface stress in both the finite-patch and the laterally-infinite limit, and the dysbiosis
contrast survives in each; the implant geometry is the native setting for the result. The eigenstrain deck is small-strain, whereas the
calibrated growth ($\lambda_z\!\sim\!1.4$) is finite; re-solving the identical thread with the
multiplicative growth split $\bF=\bF_{\!e}\bF_{\!g}$ (compressible neo-Hookean, the socket-free
\texttt{umat\_growth\_phi} UMAT under NLGEOM) reproduces the same scale-free contrasts ---
$\mathrm{DH/CH}=2.7\times$ and thread/flat $=3.6\times$, against $2.5\times$ and $4.1\times$
small-strain (Figure~\ref{fig:fem_implant}D) --- so the linearisation is adequate for this geometric
study and the conclusions are not an artefact of it.

\begin{figure}[htbp]
  \centering
  \includegraphics[width=0.80\textwidth]{fem_implant_thread.pdf}
  \caption[Implant-thread FEM: field, convergence, thread-depth sweep, finite-strain check]{Base-bonded
  biofilm patch on a smooth threaded titanium section (CPE4 plane strain; CLSM-calibrated
  volume-matched growth). \textbf{(A)}~von~Mises stress over interior thread periods: the
  growth-induced stress concentrates at the bonded thread. \textbf{(B)}~mesh convergence of the
  interior peak interface stress (the bonded-base/free-edge corner, a boundary-condition singularity,
  is excluded), bounded within $\pm7\%$ for $\ge24$ elements per thread period. \textbf{(C)}~thread-depth
  ($A/h$) sweep: a flat substrate ($A=0$) relaxes nearly stress-free, while the thread generates the
  interface stress, rising monotonically with depth. \textbf{(D)}~the scale-free $\mathrm{DH/CH}$ and
  thread/flat ratios under the small-strain eigenstrain versus the finite-strain $\bF=\bF_{\!e}\bF_{\!g}$
  UMAT agree, confirming the linearisation. The dysbiotic film adds a further
  $\mathrm{DH/CH}\approx2.5$--$2.7\times$ on top of the geometric stress.}
  \label{fig:fem_implant}
\end{figure}

The full composition$\rightarrow$stress$\rightarrow$detachment chain on the implant is
summarised intuitively in Figure~\ref{fig:fem_schematic}: the \textit{P.~gingivalis}-enriched
dysbiotic film carries a hot interface stress band over the thread and delaminates, whereas the
commensal film, with its weak and depth-spread stress, stays bonded.

\begin{figure}[htbp]
  \centering
  \includegraphics[width=0.70\textwidth]{fem_clinical_schematic.pdf}
  \caption[Mechanism schematic: why dysbiosis detaches the biofilm at the implant
  interface]{Cross-section of a titanium implant thread under the commensal (CH) and
  dysbiotic (DH) biofilm. \textit{P.~gingivalis} (dots) and the resulting growth-induced
  interface stress (red band) follow the measured depth composition: spread through the thin
  commensal film (faint interface stress, stays bonded) versus packed at the substratum in the
  thicker dysbiotic film (concentrated interface stress, delamination crack). The bar reports
  the interface von~Mises stress, $\approx\!2.5\times$ larger for the dysbiotic condition
  (implant-thread FEM interior peak, Figure~\ref{fig:fem_implant}).}
  \label{fig:fem_schematic}
\end{figure}

\paragraph{From the in-vitro contrast to a patient-level risk ranking.} The same kernel that
distinguishes the dysbiotic from the commensal film in vitro also \emph{ranks individual
patients}: carrying $J$ across the Dieckow cohort orders the ten over a more-than-twentyfold
range of detachment risk, fixed almost entirely by the
measured \textit{P.~gingivalis} abundance (the shape being a $34\times$ weaker lever). A
routine 0-D readout --- relative \textit{Pg} load from a peri-implant 16S or qPCR swab, no
depth imaging --- is therefore an adequate ordinal surrogate for the mechanical detachment
risk the full depth-resolved model would assign, turning an abundance the clinician already
measures into a ranking of the event that seeds peri-implant exposure. The claim is ordinal,
not absolute: $J$ ranks risk, while the per-patient failure threshold stays parametric
(Section~\ref{sec:fem_validation}).

\subsection{Relation to prior art, validation, and limits}
\label{sec:fem_validation}

The mechanical framework used here is closest in spirit to the finite-growth biofilm
models of B{\"o}l and co-workers, who likewise build on the multiplicative
decomposition $\bF=\bF_{\!e}\bF_{\!g}$ and even derive the resulting residual-stress
field directly from confocal biofilm geometries~\cite{Bol2009Detachment,
BoleaAlbero2014}. Two differences delimit the present contribution. First, in those
models growth is a \emph{phenomenological isotropic deposition} of new matrix, and
the attendant residual stress is continuously relaxed by a pressure-restricted,
fluid-like viscous flow; here growth is instead driven by a \emph{species-resolved
field} $\phi_\mathrm{Pg}(z)$ from the Hamilton-principle ecological model, calibrated
against CLSM depth profiles, and is anisotropic and substratum-normal. The stress we
report is therefore the \emph{composition-resolved} residual stress that a given
dysbiotic community imprints --- a quantity the deposition models, blind to who grows
where, cannot produce. Second, we deliberately model the fast, elastic limit; the
fluid-like relaxation of Bolea Albero et al.\ is the long-time complement of our
model. Incorporating it as a Prony relaxation of the elastic stress, with the shear
relaxation spectrum taken from biofilm rheometry~\cite{Shaw2004Burger}, was verified
against the analytic relaxation modulus (Figure~\ref{fig:fem_viscoelastic}; the
relaxed-to-instantaneous stress ratio matches theory to $0.0544$ vs $0.0543$, RMS
error $9\times10^{-4}$ over the full curve), confirming that the elastic result is the
short-time end of a consistent viscoelastic description.

The pressure-restricted picture of Bolea Albero et al.\ has, in addition, a
\emph{poroelastic} reading: the biofilm is a fluid-saturated porous matrix, so the
growth-induced pressure also relaxes by drainage of pore fluid through the matrix.
Modelling the growth pressurisation as an excess pore pressure in a column drained at
the biofilm--bulk surface and sealed at the substratum gives a one-dimensional Terzaghi
consolidation~\cite{Biot1941Consolidation} with drainage timescale
$\tau = H^2/c_v$, $c_v = k\,E_{\mathrm{oed}}/\gamma_w$. A finite-difference solution
matches the closed-form Terzaghi series to an RMS of $1.8\times10^{-3}$
(Figure~\ref{fig:fem_poroelastic}): the effective interface stress is half-relaxed at
$t\approx0.20\,\tau$ and essentially gone by $t\approx\tau$. The elastic residual stress
of Section~\ref{sec:fem_residual} is therefore the common short-time limit of \emph{two}
distinct relaxation channels --- viscoelastic creep of the matrix and poroelastic
drainage of the pore fluid --- both of which carry the system toward the long-time,
pressure-restricted state; which dominates is set by the ratio of the rheological to the
consolidation timescale.

\begin{figure}[htbp]
  \centering
  \includegraphics[width=0.72\textwidth]{fem_poroelastic.pdf}
  \caption[Poroelastic drainage of the growth-induced stress]{Poroelastic relaxation of
  the residual stress by pore-fluid drainage (Terzaghi consolidation).
  \textbf{(left)}~excess pore pressure dissipating from the drained biofilm--bulk surface
  toward the sealed substratum at successive times; \textbf{(right)}~the residual-stress
  fraction $1-U(t)$ versus $t/\tau$, the finite-difference solver matching the analytic
  Terzaghi series (RMS $1.8\times10^{-3}$). The elastic result is the $t\!\to\!0$ limit;
  the drained, pressure-restricted state is the $t\!\to\!\infty$ limit.}
  \label{fig:fem_poroelastic}
\end{figure}

\begin{figure}[htbp]
  \centering
  \includegraphics[width=0.56\textwidth]{fem_ve_relaxation.pdf}
  \caption[Viscoelastic stress relaxation against analytic Prony series]{Shear-stress
  relaxation of the biofilm under a held strain (two-term Prony series fitted to a
  biofilm Burgers creep model). The finite-element relaxation matches the analytic
  modulus (RMS $9\times10^{-4}$); the elastic residual stress of
  Section~\ref{sec:fem_residual} is the short-time limit of this curve.}
  \label{fig:fem_viscoelastic}
\end{figure}

\paragraph{What is and is not validated.} The growth \emph{shape} is calibrated
($\phi_\mathrm{Pg}$-driven, $r=0.88$), the elastic and viscoelastic constitutive
behaviour is verified against closed-form solutions, and all condition contrasts
(residual stress, delamination) are reported at identical parameters so that only the
composition differs. What remains parametric is the absolute stress \emph{scale}: the
biofilm modulus spans roughly \emph{seven orders of magnitude} --- from
$E\sim\mathrm{O}(1)\,\mathrm{MPa}$ by short-time AFM nanoindentation of oral
biofilm~\cite{Cense2021OralBiofilm}, through $G'\sim\mathrm{O}(10)\,\mathrm{kPa}$ by
oscillatory shear rheology~\cite{BiofilmMatrixMechanics2025}, down to
$G\sim\mathrm{O}(1\text{--}10)\,\mathrm{Pa}$ by long-time creep
rheometry~\cite{Stoodley2002Material}, with no depth-resolved oral-biofilm modulus or
residual-stress measurement available --- which is why the result is reported
modulus-agnostically as $\sigma/\mu$, calibrated in shape but parametric in scale
($\gamma$), with depth-resolved AFM nanoindentation or shear-flow detachment onset versus
\textit{P.~gingivalis} load identified as the next experiment.

\paragraph{Per-field heterogeneity and the level of the claim.} The depth shape that drives
the eigenstrain is pooled across confocal fields of view. Resolving it instead per individual
field --- all $84$ titanium fields, recovered by the same colocalisation decode --- shows that
the detachment driver $J$ does \emph{not} separate the dysbiotic from the commensal condition
field-by-field (Mann--Whitney $p\approx0.8$; Figure~\ref{fig:fem_perfov}), and that the
dysbiotic fields are the more heterogeneous (coefficient of variation $0.56$ against $0.34$).
The condition contrast is therefore a property of the \emph{pooled} profile and its time course
--- the absolute, base-localised \textit{P.~gingivalis} that accumulates over the dysbiotic
series --- and not of any single field's composition, which is too variable (and, for the
dysbiotic arm, too sparsely sampled: $2$--$3$ fields at the later days) to classify
mechanically. The dysbiosis$\rightarrow$detachment claim is accordingly made at the population
and trajectory level; a per-field mechanical risk map is explicitly \emph{not} supported by the
present data, and the pooled depth shape used throughout is the appropriate level of resolution
for it.

\begin{figure}[htbp]
  \centering
  \includegraphics[width=0.72\textwidth]{fem_perfov_heterogeneity.pdf}
  \caption[Per-field-of-view detachment driver: no single-field separation]{\textbf{(left)}
  Depth-resolved \textit{P.~gingivalis} composition $\phi_\mathrm{Pg}(z)$ for all $84$ confocal
  fields (faint) with per-condition means; \textbf{(right)} the per-field detachment driver $J$,
  which does not separate the conditions (Mann--Whitney $p\approx0.8$) and is markedly more
  heterogeneous for the dysbiotic arm. The mechanical contrast lives in the pooled,
  absolute-magnitude profile and its time course, not in single-field composition.}
  \label{fig:fem_perfov}
\end{figure}

\paragraph{Two real datasets, two layers, and the strength of each.} The framework is driven
by real data at two layers that must not be conflated: the patient cohort (Dieckow~\textit{et
al.}\ peri-implant 16S~\cite{dieckow2024}) is real but composition-only and drives the
ecological model of Chapters~\ref{ch:heine}--\ref{ch:dieckow}, while the depth profile
$\phi_\mathrm{Pg}(z)$ that sets the eigenstrain is \emph{borrowed} from the in-vitro HOBIC
confocal imaging~\cite{Heine2025PeriImplant} (colocalisation decode of
Section~\ref{sec:fem_growth}), so the two join as a measured patient \emph{load} $\times$ an
in-vitro depth \emph{shape}. The borrowed shape is a weak lever and the measured load the strong
one: projecting the cohort through the boundary-layer driver
$J=\langle e^{-z/z_0}\phi_\mathrm{Pg}(z)\rangle_z$ ($z_0=0.25$, Section~\ref{sec:fem_residual})
gives a large detachment-risk range set entirely by the measured Pg load (CV $0.86$,
$J\in[0,\,0.028]$), whereas substituting the entire commensal shape for the dysbiotic one at
fixed load changes $J$ by only $8\%$ --- the inter-patient spread is thus $34\times$ wider than
the shape-uncertainty band, so per-patient risk is governed by the 0-D load the clinic can
measure. This low shape-leverage holds for the driver $J$ that ranks risk, not for the absolute
stress \emph{scale} (parametric in the biofilm modulus) nor for the near-substratum depth
\emph{kernel} on which the real-tooth run (Section~\ref{sec:fem_tooth}) still relies; removing
the borrowing needs depth-resolved imaging of a patient biofilm and per-patient detachment onset
versus Pg load to test the linear $J$--load law.

\paragraph{Surface wrinkling: a second mechanical signature of dysbiosis.}
Delamination is the failure of the film--substrate interface; the growth-induced
compression also destabilises the \emph{free surface}. A stiff matrix crust growing on a
softer, non-growing foundation (a two-material film, the configuration of a mature
biofilm with a denser surface layer) is forced into in-plane compression and wrinkles.
Carrying this to elevated growth, the system stiffness matrix develops negative
eigenvalues at a critical growth --- a genuine bifurcation, not merely an amplified
imperfection --- and the post-bifurcation surface selects a wavelength that is
\emph{condition-specific}: the dysbiotic crust crosses the threshold into a short-wavelength
wrinkle, whereas the commensal crust, growing less, remains in the long-wavelength
global-bending mode and does not wrinkle (Figure~\ref{fig:fem_wrinkling}). Surface
wrinkling thus joins interface delamination as a mechanical instability that the
dysbiotic composition triggers and the commensal one does not, both driven by the same
$\phi_\mathrm{Pg}$-set growth.

\begin{figure}[htbp]
  \centering
  \includegraphics[width=0.72\textwidth]{fem_wrinkling.pdf}
  \caption[Growth-induced surface wrinkling: dysbiotic vs commensal]{Growth-induced
  wrinkling of a stiff crust on a soft foundation. \textbf{(left)}~surface-deflection
  amplitude versus growth: a super-linear onset marks the instability;
  \textbf{(right)}~the final surface profile --- the dysbiotic crust selects a short
  wrinkle wavelength ($\lambda\!\approx\!1.8$), while the commensal crust stays in the
  long-wavelength bending mode. Identical material and imperfection; only
  $\phi_\mathrm{Pg}$ differs.}
  \label{fig:fem_wrinkling}
\end{figure}

\paragraph{Outlook.} Two further extensions are left for future work. The constitutive
routine ports to a second commercial solver essentially unchanged via the
predefined-field path of Section~\ref{sec:fem_method}, enabling its integration into an
established implant-mechanics framework. And the crude neo-Hookean elastic factor can be
replaced by the worm-like-chain network model of EPS
mechanics~\cite{EhretBol2013} for a more physical constitutive law that captures the
strain-stiffening the present model omits.

\paragraph{Implant-thread extensions.} Beyond the threaded result of
Section~\ref{sec:fem_delamination}, a family of finite-element extensions on the implant
geometry (Figure~\ref{fig:fem_implant_ext}) probes robustness. A superimposed micro-thread
concentrates the interior interface stress further (peak $\sigma_\mathrm{vM}\!\sim\!1030$ against
$\sim\!610$ for the plain thread), and an AT2 phase-field damage field driven by the measured
interface-stress profile delaminates the dysbiotic interface $2.3\times$ more than the commensal.
The time-dependent interface state is coupling-dependent: the biofilm Prony series relaxes the
thread stress \emph{viscoelastically} to $\sim\!10\%$ of its instantaneous value over $t\gg\tau$,
whereas a coupled \emph{poroelastic} model (CPE4P) with growth carried as a skeleton eigenstrain
gives the opposite sign --- drainage \emph{raises} the effective interface stress ($\approx\!1.55\times$,
toward the drained elastic value), opposite to the Terzaghi picture of
Section~\ref{sec:fem_validation}, so the long-time state depends on which growth--fluid coupling
dominates. Curvature adds a three-dimensional effect a plane-strain section omits: an axisymmetric
screw carries an inner-interface hoop stress scaling with $1/R_0$, and a full three-dimensional
helical model (C3D8; Figure~\ref{fig:fem_implant_3d}) shows the growth-induced stress wrapping the
screw as a helical band that, under occlusal loading, acquires a cyclic amplitude
($\Delta\sigma_\mathrm{vM}\!\sim\!230$ at the dysbiotic interface) driving interface fatigue over
$\sim\!10^6$ chewing cycles, fastest at the dysbiotic interface ($\approx\!3\times$ the commensal mean),
with the outer flow-shear traction adding only $\sim\!2\%$. A two-way chemo--mechanical coupling through
a stress-dependent permeability is the natural closure; the production continuation is a user-material
implementation in the institute's ANSYS implant-mechanics framework.

\paragraph{Generic-implant design study.} Replacing the patient tooth-23 root with a generic standard
screw ($\varnothing 4.1$\,mm, $1.0$\,mm pitch) osseointegrated in the real two-layer mandible and loaded
at the ISO~14801 $30^\circ$ oblique condition reproduces the same mechanics while adding a quantitative
design map (Figure~\ref{fig:fem_implant_design}). Peri-implant stress is governed first by diameter:
widening the screw from $\varnothing 3.5$ to $\varnothing 4.8$\,mm lowers the thread-root peak von~Mises
stress from $171$ to $60$\,MPa and raises the coupon stiffness $3.5\times$, whereas length
($8$--$12$\,mm) and pitch ($0.6$--$1.0$\,mm) are secondary; the load \emph{angle} dominates, a $30^\circ$
oblique bite raising thread-root stress $\approx\!9\times$ and crestal-bone stress $\approx\!7\times$ over
axial load. Resolving the concentration requires quadratic elements (linear C3D4 under-predict by
$\sim\!10\%$, so C3D10 are used); all peak titanium stresses ($60$--$230$\,MPa) stay within the
Ti-6Al-4V yield. This connects the detachment mechanism to a design guideline --- wide diameter,
near-axial loading --- and feeds the ANSYS framework above.

\paragraph{From mechanics to disease dynamics: bistability.} The mechanical results feed a host-response
cascade that closes the loop of this thesis: dysbiotic drive raises inflammation, inflammation
up-regulates RANKL/OPG and osteoclastic bone loss, and the resulting recession concentrates crestal
stress that --- only under inflammation --- aggravates resorption, while inflammation also feeds the
anaerobes back (heme-iron for \emph{P.\,gingivalis}), closing an ecology--host positive-feedback loop.
With this loop the system becomes \emph{bistable} (Figure~\ref{fig:fem_bistability}): healthy and
diseased attractors coexist across a threshold, so a transient insult can lock the implant into
progressive loss with hysteresis, reframing peri-implantitis as a tipping-point disease and giving a
mechanistic basis for the primacy of prevention. Companion analyses (a reversibility-window map, an
anti-RANKL therapy-node comparison, a literature-anchored bone-loss rate $k_L=0.022\,\mathrm{wk}^{-1}$,
and a design-buys-time coupling) are reported in the project record; all are ODE/post-processing models,
and a per-timepoint Bayesian calibration is the natural next step once a longitudinal cohort pairing
microbiome with radiographic bone loss (e.g.\ ENA~PRJNA1215005) becomes available.

\begin{figure}[htbp]
  \centering
  \includegraphics[width=0.78\textwidth]{fem_implant_extensions.pdf}
  \caption[Implant-thread finite-element extensions]{Extensions on the threaded geometry.
  \textbf{(A)}~a superimposed micro-thread (two-scale roughness) concentrates the interface stress.
  \textbf{(B)}~an AT2 phase-field damage field driven by the measured thread interface-stress profile:
  the dysbiotic interface accrues $2.3\times$ more damage than the commensal. \textbf{(C)}~with the
  biofilm Prony series the thread interface von~Mises stress relaxes to $\sim\!10\%$ of its
  instantaneous value over $t\gg\tau$.}
  \label{fig:fem_implant_ext}
\end{figure}

\begin{figure}[htbp]
  \centering
  \includegraphics[width=0.76\textwidth]{fem_implant_3d.pdf}
  \caption[Three-dimensional helical implant-screw FEM]{Full 3-D helical implant-screw biofilm
  (C3D8, isotropic volume-matched growth). \textbf{(A)}~the bonded interface coloured by von~Mises
  stress; \textbf{(B)}~the same interface unrolled in (azimuth, axial height), where the helical thread
  appears as diagonal stress bands. The growth-induced interface stress follows the thread helix; it
  is the substrate for the occlusal-fatigue and flow-shear load cases discussed in the text.}
  \label{fig:fem_implant_3d}
\end{figure}

\begin{figure}[htbp]
  \centering
  \includegraphics[width=0.78\textwidth]{fem_implant_design.pdf}
  \caption[Generic-implant design study (ISO 14801 coupon)]{Generic screw-implant design map from an
  ISO~14801-style coupon (titanium screw osseointegrated in a clamped bone holder, loaded at $30^\circ$).
  \textbf{(A)}~implant \emph{diameter} dominates: thread-root von~Mises stress falls and coupon stiffness
  rises sharply with $\varnothing$. \textbf{(B)}~length, pitch and taper are secondary (bars near the
  baseline, dashed). \textbf{(C)}~load \emph{angle} dominates over geometry --- a $30^\circ$ oblique bite
  multiplies the thread-root stress, crestal-bone stress and displacement by $\approx\!9$, $7$ and $27$
  over axial loading (log scale). \textbf{(D)}~quadratic C3D10 elements capture the thread-root
  concentration that linear C3D4 under-predict.}
  \label{fig:fem_implant_design}
\end{figure}

\begin{figure}[htbp]
  \centering
  \includegraphics[width=0.80\textwidth]{fem_periimplantitis_bistability.pdf}
  \caption[Bistability and hysteresis of the peri-implantitis cascade]{Closing the ecology--host
  feedback loop makes the disease cascade bistable. \textbf{(A)}~phase plane of dysbiotic load $B$ versus
  inflammation $I$ with two stable fixed points (healthy / diseased). \textbf{(B)}~the RANKL/OPG ratio
  traced as the plaque burden is raised then lowered describes a hysteresis loop --- the diseased branch
  persists well below the burden that triggered it (prevention $\ll$ cure). \textbf{(C)}~a transient
  insult tips the patient permanently into the diseased basin.}
  \label{fig:fem_bistability}
\end{figure}

\subsection{Generalisation to patient-specific geometry}
\label{sec:fem_tooth}

\paragraph{The substratum is, in vivo, a titanium implant.} This work's clinical setting, and
the Dieckow cohort's, is peri-implant: the biofilm colonises the transmucosal abutment in the
peri-implant sulcus~\cite{dieckow2024}, and the substratum used throughout (the Ti-mounted
in-vitro biofilm of Section~\ref{sec:fem_growth}) is already titanium. The implant is thus
solved directly --- the threaded-titanium FEM of Section~\ref{sec:fem_delamination}
(Figure~\ref{fig:fem_implant}; mechanism in Figure~\ref{fig:fem_schematic}) is the
\emph{primary} geometric result. The conformal real-tooth run that follows is therefore only a
smooth-geometry robustness control on an independent curved anatomy, not a claim that the tooth
is the operative surface.

The idealised column isolates the substratum-interface mechanism but abstracts away the
curved, patient-specific anatomy on which a biofilm actually grows. To establish that the
result is not an artefact of the simplified geometry, the \emph{identical}
composition-driven growth model was applied to a conformal biofilm mesh on a \emph{real
human tooth} (tooth~23 of the open-access Open-Full-Jaw
dataset~\cite{gholamalizadeh2022open}): a thin biofilm of $86\,160$ tetrahedra bonded at
its inner surface to the tooth (the substratum) and free at its outer surface
(Figure~\ref{fig:fem_tooth}). Four measures make the transfer examiner-robust. First, to
avoid the volumetric locking of linear tetrahedra near incompressibility, \emph{hybrid}
elements (C3D4H) are used at the same Poisson ratio $\nu=0.45$ as the column; the
field-driven user material and its consistent tangent carry over unchanged, converging in
one to two Newton iterations per increment, and a single-tetrahedron free-swelling test
reproduces the prescribed growth at machine-zero stress ($J_g=1.20$,
$\lvert\sigma\rvert\sim10^{-9}\mu$), confirming the implementation on the tetrahedral
element. Second, the conformal mesh is \emph{mesh-converged} on the real geometry:
refining the through-thickness resolution $N_\ell=4\!\to\!8\!\to\!16$ changes the peak
stress by only $-0.8\%$ then $-0.2\%$ (Figure~\ref{fig:fem_tooth}, right), so the
eight-layer mesh is converged. Third, the growth is driven by the \emph{measured DH depth
profile} $\phi_\mathrm{Pg}(z)$ (Section~\ref{sec:fem_growth}) mapped through the biofilm
thickness, not a uniform value. Fourth, the stress is reported as $\sigma/\mu$; the
von~Mises invariant is used because the curved surface has no global in-plane axis, but
the substratum layer is in fact \emph{uniformly compressive}
($\sigma_{\min}\approx-0.57\,\mu$ at every interface element), consistent with the
compressive in-plane stress of the column (Section~\ref{sec:fem_residual}).

With these controls the composition-resolved coupling survives the transfer to real
anatomy. The reading is deliberately honest: because the growth is driven by
$\phi_\mathrm{Pg}$, a correlation between stress and composition is \emph{built into} the
model; what the real-geometry run tests is whether the curvature, the thin conformal layer
and the tie to the tooth \emph{destroy} it --- and they do not. Through the biofilm
thickness the residual stress tracks the local \textit{P.~gingivalis} fraction with a
per-element Pearson correlation $r=0.92$ (Spearman $0.90$). Running the \emph{commensal}
profile through the identical model then recovers the clinical contrast on the real tooth:
with only the composition changed, the dysbiotic film develops a far more
\emph{stratified} stress field than the commensal one --- $3.8\times$ the through-depth
amplitude and a $1.4\times$ higher peak (Figure~\ref{fig:fem_tooth}, centre), the
heterogeneity that drives interfacial failure. The
composition\,$\rightarrow$\,growth\,$\rightarrow$\,residual-stress chain is thus
geometry-portable: with the same verified material model and a public clinical mesh it
reproduces, on patient-specific anatomy, both the composition tracking and the
dysbiotic-versus-commensal contrast of the idealised column.

\begin{figure}[htbp]
  \centering
  \includegraphics[width=0.72\textwidth]{fem_tooth_geometry.pdf}
  \caption[Composition-resolved residual stress on a real tooth]{\textbf{(left)} The
  depth-graded DH growth model on a conformal biofilm over a real tooth (Open-Full-Jaw
  tooth~23; hybrid C3D4H, $\nu=0.45$), coloured by residual von~Mises stress
  $\sigma_\mathrm{vM}/\mu$. \textbf{(centre)} Through-thickness profiles: with only the
  composition changed, the dysbiotic film is $3.8\times$ more stratified than the
  commensal one (per-element $\sigma$--$\phi_\mathrm{Pg}$ correlation $r=0.92$ for DH).
  \textbf{(right)} The result is mesh-converged on the real geometry (peak stress
  $N_\ell=8\to16$: $-0.2\%$). This is a \emph{smooth-geometry robustness control} for the
  threaded-implant result (Figure~\ref{fig:fem_implant}): the composition-resolved mechanism and
  the dysbiotic-versus-commensal contrast of Section~\ref{sec:fem_residual} survive transfer to
  an independent, curved patient anatomy.}
  \label{fig:fem_tooth}
\end{figure}

\paragraph{A full peri-implant assembly: implant, tooth, and shared bone.} The single-column
and single-tooth runs isolate the substratum-interface mechanism; the clinical reality is a
titanium implant standing \emph{next to} natural teeth in a shared alveolar bone, which is the
setting peri-implant disease actually occupies. To close that gap the same composition-driven
growth model was embedded in a coupled multi-material assembly --- a threaded titanium screw
implant at the (extracted) tooth-23 site, the neighbouring tooth-24 with its periodontal ligament, a
conforming dysbiotic biofilm collar, and the shared mandibular bone (real Open-Full-Jaw
anatomy)~\cite{gholamalizadeh2022open} --- and loaded occlusally
(Figure~\ref{fig:fem_tier2b}). The assembly reproduces the textbook contrast that motivates
the whole chapter: load transmitted through the shared bone concentrates at the
\emph{crestal} peri-implant bone around the osseointegrated, ligament-free implant, whereas the
neighbouring tooth distributes it along the periodontal ligament down the root --- the
mechanical signature of peri-implant marginal bone loss. The composition-resolved biofilm
mechanism of this chapter thus sits, without modification, inside a realistic implant--tooth--bone
model: the geometry that defines implant dentistry.

\begin{figure}[htbp]
  \centering
  \includegraphics[width=0.78\textwidth]{fem_implant_screw3d.pdf}
  \caption[Coupled implant--tooth--bone peri-implant assembly]{The composition-driven model in
  a full peri-implant assembly, rendered from the real tetrahedral mesh (51k nodes, 252k C3D4)
  as shaded boundary surfaces with the front alveolar bone clipped away.
  \textbf{(A)}~the multi-material geometry: a threaded titanium screw implant (tooth-23 site),
  the neighbouring tooth-24 with periodontal ligament and dentin, dysbiotic biofilm collars, in
  shared alveolar bone. \textbf{(B)}~occlusal-load von~Mises stress: load through the shared bone
  concentrates at the \emph{crestal} bone around the ligament-free implant while the tooth
  distributes it along the root --- the peri-implant marginal-bone-loss signature. The
  composition-resolved biofilm mechanism transfers unchanged into this realistic setting.}
  \label{fig:fem_tier2b}
\end{figure}

\section{Clinical Applications}
\label{sec:integration_clinical}

The modelling framework has three concrete clinical translational directions, each
grounded in results from the preceding chapters.

\paragraph{Early dysbiosis detection.}
The guild-level red/green dysbiosis ratio derived from the Dieckow interaction matrix
separates peri-implantitis from health/mucositis on 127 independent cross-sectional
samples (Spearman $\rho=0.46$, $p=5\times10^{-8}$; Chapter~\ref{ch:dieckow}).
Applied to routine 16S amplicon data from implant abutment plaque, the ratio requires
no model fitting --- only taxonomic classification to the ten guild taxonomy
(Table~\ref{tab:dieckow_guilds}) followed by a log-ratio calculation.
This makes it deployable as a single-number peri-implant risk score, with the
current threshold ($\mathrm{R/G}=0$) distinguishing healthy and mucositic from
peri-implantitis patients ($p<10^{-4}$). Extending the Joshi validation to
longitudinal series from implant maintenance visits would enable prospective
monitoring of trajectory toward or away from dysbiosis.

\paragraph{Treatment optimisation via growth-vector personalisation.}
The Hamilton model separates shared ecology ($\bA$, community-level interaction
structure) from patient-specific growth rates ($\bb^{(p)}$). Since $\bb^{(p)}$ is
identifiable from a single W1 sample, a personalised forward simulation is
available at the first clinical timepoint. In principle, any intervention that
shifts $b_i^{(p)}$ --- antibiotic suppression of late colonisers, probiotic
supplementation of early colonisers, or mechanical debridement that resets the
initial condition --- can be evaluated \emph{in silico} before application. The CT1/CT2
clustering of growth vectors (Figure~\ref{fig:dieckow_bpca} in
Appendix~\ref{app:dieckow_supp}) shows that this personalisation is clinically
meaningful: the two community types differ systematically in which guilds carry
strong positive $b_i$.

\paragraph{Probiotic and microbiome-steering design.}
The UMAP analysis of the in-vitro posteriors (Chapter~\ref{ch:heine}) identifies the
commensal CS/CH attractors as the target region of the interaction-matrix landscape.
The depth-ordering result confirms that the $\textit{So}\to\textit{Vei}$ cross-feeding
edge (Chapter~\ref{ch:integration}) is spatially encoded and therefore structure-dependent:
steering toward the commensal attractor requires not just population-level abundance
shifts but restoration of the stratified biofilm architecture. A probiotic strategy
that enriches early colonisers (\textit{S.~oralis}, \textit{Actinomyces}) to shift
the community away from the Fn--Pg co-aggregation block is consistent with both the
ODE attractor landscape and the FISH depth-ordering data.

\section{Outlook: A Unified Gradient-Flow Perspective}
\label{sec:integration_outlook}

The three layers of this chapter --- Bayesian interaction inference, the
reaction--diffusion spatial extension, and the finite-element mechanical realisation
--- are not disjoint models but three readings of one variational structure. The
Extended Hamilton Principle of Chapter~\ref{ch:theory}, applied to an overdamped
internal-variable field with a quadratic dissipation potential, yields by stationarity
an $L^2$ gradient flow of the free energy --- the mathematical core of the phase-field
method~\cite{AllenCahn1979} --- whose monotone Lyapunov decay is exactly the
thermodynamic consistency of Section~\ref{sec:theory_hamilton}. Which model results is
selected only by the metric and by whether the field is conserved:
\begin{center}\small
\begin{tabular}{@{}lll@{}}
\toprule
metric / conservation & resulting equation & instance in this thesis\\
\midrule
Euclidean, non-conserved & Allen--Cahn          & damage / fracture field\\
conserved flux           & Cahn--Hilliard       & biomass demixing with depth\\
Fisher--Shahshahani      & replicator           & inference of Ch.~\ref{ch:heine}--\ref{ch:dieckow}\\
\bottomrule
\end{tabular}
\end{center}
Closing the spatial model in conserved form recovers a reactive Cahn--Hilliard
description~\cite{CahnHilliard1958} in which the vertical \textit{P.~gingivalis}
sequestration of FISH (Section~\ref{sec:pde_fish}) is a spinodal-like demixing into
depth niches; regularising the cohesive delamination of
Section~\ref{sec:fem_delamination} with a damage field recovers AT2 phase-field
fracture~\cite{Bourdin2008Variational,Miehe2010PhaseField} driven by the same growth
energy. Growth, fracture, and ecological inference are therefore one Extended-Hamilton
gradient flow under different metrics and interface terms, and this phase-field
completion --- not extrapolation but continuation --- is the principal methodological
direction beyond this thesis.
